% ---------------------------------------------------------------------
% --- CAPÍTULO 5 - CONCLUSÕES
% --- Conteúdo extraído VERBATIM de TCC_PARKON_V3.pdf (página 31),
% --- fonte da verdade. Fragmento LaTeX (sem \documentclass/preâmbulo),
% --- incluído via \include em tese.tex. Arquivo em UTF-8.
% ---
% --- Título de origem (corpo do PDF): "CONCLUSÕES" (capítulo 5).
% ---   A classe PPGC_UFF formata o título do capítulo em maiúsculas;
% ---   por isso usa-se \chapter{Conclusões} (mesma saída visível),
% ---   mantendo a coerência com cap1.tex/cap2.tex (Requisitos 1.2 e 1.6).
% ---   \label{cap:conclusao} preservado conforme convenção do projeto
% ---   (referenciado por \autoref em cap1.tex).
% ---
% --- MAPEAMENTO DE CITAÇÕES (fonte usa numeração [n]; convertidas para \cite):
% ---   [2]  -> \cite{flutter2025}    "Application development using flutter" (sem autor; @misc/online)
% ---   [10] -> \cite{goldstein2025}  GOLDSTEIN, S. "Criação de plataforma de geocoding baseada em serviços Google Maps"
% ---   TODO: [CITATION] As referências da fonte NÃO trazem ano de publicação
% ---         explícito (apenas "Acesso em 08 de dez. 2025"). A chave
% ---         flutter2025 é PROVISÓRIA (entrada [2] não possui autor; padrão
% ---         {sobrenomePrimeiroAutor}{ano} não se aplica diretamente).
% ---         Chaves/anos serão criados e reconciliados nas tarefas 7.1 e 7.2.
% ---
% --- OBSERVAÇÕES DE MIGRAÇÃO:
% --- - Não há citações diretas com mais de três linhas (nenhum bloco \quoting é necessário).
% --- - Não há figuras nem tabelas neste capítulo.
% --- - Não há referências cruzadas a figuras/tabelas/seções (nenhum \autoref é necessário).
% --- - Nenhum acrônimo definido (uff, api, apk, cli, rf, rnf) ocorre no texto deste capítulo.
% --- - O conteúdo de "trabalhos futuros" está integrado ao último parágrafo
% ---   ("Como perspectivas futuras..."); a fonte NÃO possui uma seção separada,
% ---   portanto não se cria \section (preserva a estrutura VERBATIM, Req. 1.1/4.2).
% ---------------------------------------------------------------------

\chapter{Conclusões}
\label{cap:conclusao}

O desenvolvimento da plataforma Parkon demonstrou o potencial das soluções tecnológicas para otimizar o processo de busca, reserva e avaliação de vagas de estacionamento em ambientes urbanos. Utilizando o framework Flutter~\cite{flutter2025} para o aplicativo mobile, Go para o backend e React com Next.js para o portal web administrativo, foi possível criar um ecossistema flexível, eficiente e escalável, capaz de integrar funcionalidades essenciais como geolocalização em tempo real, reserva antecipada, avaliação colaborativa de vagas e gestão completa de estacionamentos.

Um dos principais desafios enfrentados foi a integração do serviço de geolocalização do Google~\cite{goldstein2025} com o sistema de cache baseado em MongoDB, garantindo desempenho e disponibilidade das informações mesmo em situações de instabilidade de conexão. A adoção de tecnologias abertas e padrões amplamente utilizados permitiu maior interoperabilidade e facilidade de manutenção, além de ampliar as possibilidades de crescimento e inovação da plataforma.

A proposta colaborativa do Parkon, viabilizada pelo portal web administrativo que permite a qualquer estacionamento se cadastrar e gerenciar suas vagas de forma autônoma, representa um avanço em relação às soluções tradicionais, tornando o serviço mais democrático e abrangente. Os resultados obtidos indicam benefícios significativos para os usuários e para a mobilidade urbana, promovendo eficiência, praticidade e sustentabilidade.

Como perspectivas futuras, sugere-se a implementação de dashboards avançados com relatórios de ocupação, integração com sistemas de pagamento online, notificações push para confirmação de reservas e ampliação da base de estacionamentos parceiros. O trabalho realizado contribui para o avanço da pesquisa e do desenvolvimento de soluções inovadoras na área de mobilidade urbana, evidenciando a importância da tecnologia como ferramenta para transformar o cotidiano das cidades.
